%! Permutations and Transposes — Unit 2, Lesson 2.4 — Guided Notes (Answer Key)
\documentclass[10pt]{article}
\usepackage{linalg-article}
\usepackage{linalg-key}   % pulls in -boxes; do NOT also load -boxes
\newcommand{\TallMath}[1]{$\displaystyle #1\rule[-1.4em]{0pt}{3.2em}$}
\newcommand{\xx}{\mathbf{x}}
\newcommand{\bb}{\mathbf{b}}
\newcommand{\cc}{\mathbf{c}}
\newcommand{\PP}{P}
\newcommand{\tp}{\mathsf{T}}

\begin{document}

\pageheader{Unit 2, Lesson 2.4}{Guided Notes --- Answer Key}

\vspace{0.12in}

\begin{objectivebox}
\textbf{By the end of this lesson, I will be able to\dots}
\begin{itemize}\setlength{\itemsep}{2pt}
  \item build a \emph{permutation matrix} $\PP$ and use it to swap the rows of a matrix;
  \item fix a \emph{zero pivot} with a row swap and write the factorization $\PP A=LU$;
  \item \emph{transpose} a matrix and recognize \emph{symmetric} matrices ($A^{\tp}=A$);
  \item explain why swapping two rows never changes the solution to $A\xx=\bb$.
\end{itemize}
\end{objectivebox}

\vspace{0.10in}

\begin{vocabbox}
\small
Fill in each term as we build it together.
\par\vspace{2pt}
\vocabans{Permutation matrix $\PP$}{the identity with its rows reordered; multiplying $\PP A$ carries out that same row reordering on $A$}
\vocabans{Zero pivot / row exchange}{a $0$ where a pivot is needed; swap in a lower row (a row exchange) to bring up a nonzero pivot}
\vocabans{Factorization $\PP A=LU$}{swap rows first ($\PP A$), then factor into lower$\times$upper as usual --- the fix when $A=LU$ is blocked}
\vocabans{Transpose $A^{\tp}$}{flip $A$ across its diagonal; row $i$ becomes column $i$, so entry $(i,j)$ moves to $(j,i)$}
\vocabans{Symmetric matrix}{a square matrix with $A^{\tp}=A$ --- its own mirror image across the diagonal}
\end{vocabbox}

\begin{hookbox}
Try to start elimination on the matrix below, the way we did all last lesson:
\[
  A = \begin{bmatrix} 0 & 2 \\ 1 & 3 \end{bmatrix}.
\]
\begin{itemize}\setlength{\itemsep}{1pt}
  \item The first pivot is the top-left entry --- but it is $0$. Elimination can't start. What simple
        move rescues the system?
        \ansline{Swap the two rows (equations): then the top-left entry is $1$, a good pivot, and elimination proceeds.}
\end{itemize}
\end{hookbox}

\begin{notesbox}{1. A Permutation Matrix Swaps Rows}
A \textbf{permutation matrix} $\PP$ is just the identity $I$ with its rows reordered. The $2\times2$
\emph{swap} matrix is
\[
  \PP = \begin{bmatrix} 0 & 1 \\ 1 & 0 \end{bmatrix},
  \qquad
  \PP\begin{bmatrix} a \\ b \end{bmatrix} = \begin{bmatrix}{\color{keyred}\mathbf{b}}\\[2pt]{\color{keyred}\mathbf{a}}\end{bmatrix}.
\]
Multiplying on the left, $\PP A$ swaps the two \emph{rows} of $A$:
\[
  \PP\begin{bmatrix} 0 & 2 \\ 1 & 3 \end{bmatrix}
  = \begin{bmatrix}{\color{keyred}\mathbf{1}} & {\color{keyred}\mathbf{3}}\\[2pt]{\color{keyred}\mathbf{0}} & {\color{keyred}\mathbf{2}}\end{bmatrix}.
\]
$\PP$ never changes the numbers --- it only changes their \emph{order}. And swapping twice puts
everything back, so $\PP\PP=I$; a permutation \emph{undoes itself}.
\end{notesbox}

\begin{notesbox}{2. Fixing a Zero Pivot: $\PP A=LU$}
Back to the hook, $A=\begin{bsmallmatrix} 0&2\\1&3 \end{bsmallmatrix}$. The zero pivot blocks elimination,
so \textbf{swap the rows first} with $\PP=\begin{bsmallmatrix} 0&1\\1&0 \end{bsmallmatrix}$:
\[
  \PP A = \begin{bmatrix} 1 & 3 \\ 0 & 2 \end{bmatrix}.
\]
Now the pivots are ${\color{keyred}\mathbf{1}}$ and ${\color{keyred}\mathbf{2}}$ --- both nonzero, so
elimination runs. Here $\PP A$ is \emph{already} upper triangular, so
$U=\begin{bsmallmatrix} 1&3\\0&2 \end{bsmallmatrix}$ and no elimination step is needed: $L=I$. The
factorization is
\[
  \PP A = LU, \qquad L=\begin{bmatrix} 1 & 0 \\ 0 & 1 \end{bmatrix},\quad U=\begin{bmatrix} 1 & 3 \\ 0 & 2 \end{bmatrix}.
\]
\textbf{The rule:} whenever a zero pivot appears, a permutation $\PP$ swaps rows to bring up a nonzero
pivot, and then $\PP A$ factors the ordinary way. For bigger matrices, after the swap you run 2.3's
elimination and the multipliers fill $L$ as usual (you'll do a $3\times3$ in the activity).
\end{notesbox}

\begin{notesbox}{3. Why the Swap is Safe: Solving $A\xx=\bb$}
Swapping two rows just \emph{reorders the equations} --- so the solution can't change, as long as we swap
the same two entries of $\bb$. To solve $A\xx=\bb$ we solve $\PP A\xx=\PP\bb$ instead.

\vspace{3pt}
Take $A=\begin{bsmallmatrix} 0&2\\1&3 \end{bsmallmatrix}$ and $\bb=\begin{bsmallmatrix} 2\\4 \end{bsmallmatrix}$.
Swap both: $\PP A=\begin{bsmallmatrix} 1&3\\0&2 \end{bsmallmatrix}$ and
$\PP\bb=\begin{bsmallmatrix} 4\\2 \end{bsmallmatrix}$. Since $L=I$, we back-solve $U\xx=\PP\bb$:
\[
  \begin{bmatrix} 1 & 3 \\ 0 & 2 \end{bmatrix}\begin{bmatrix} x \\ y \end{bmatrix} = \begin{bmatrix} 4 \\ 2 \end{bmatrix}
  \;\Rightarrow\; 2y=2 \Rightarrow y = {\color{keyred}\mathbf{1}},\quad x+3y=4 \Rightarrow x = {\color{keyred}\mathbf{1}}.
\]
So $(x,y)=(\,\ans{$1,\,1$}\,)$. Check it in the \emph{original} order: $0(1)+2(1)=2$ \checkmark\ and
$1(1)+3(1)=4$ \checkmark. Same equations, just reordered --- same answer.
\end{notesbox}

\begin{notesbox}{4. The Transpose and Symmetric Matrices}
The \textbf{transpose} $A^{\tp}$ flips $A$ across its diagonal: \emph{row $i$ becomes column $i$}, so the
entry at $(i,j)$ moves to $(j,i)$. Rows turn into columns:
\[
  A = \begin{bmatrix} 1 & 2 & 0 \\ 3 & 1 & 4 \end{bmatrix}
  \;\Rightarrow\;
  A^{\tp} = \begin{bmatrix}{\color{keyred}\mathbf{1}} & {\color{keyred}\mathbf{3}}\\[2pt] {\color{keyred}\mathbf{2}} & {\color{keyred}\mathbf{1}}\\[2pt] {\color{keyred}\mathbf{0}} & {\color{keyred}\mathbf{4}}\end{bmatrix}.
\]
A square matrix is \textbf{symmetric} when $A^{\tp}=A$ --- it is its own mirror image across the
diagonal. Circle the symmetric one:
\[
  S = \begin{bmatrix} 2 & 5 \\ 5 & 3 \end{bmatrix}, \qquad
  T = \begin{bmatrix} 1 & 2 \\ 3 & 4 \end{bmatrix}.
  \qquad\text{Symmetric: } \ans{$S$ (its $5$s match; $T$ has $2\ne3$)}
\]

\vspace{2pt}
Two rules we'll reuse. Transpose \textbf{reverses the order} of a product (2.2's socks-and-shoes again):
\[
  (AB)^{\tp} = {\color{keyred}\mathbf{B^{\tp}}}\,{\color{keyred}\mathbf{A^{\tp}}}\quad\text{(the reversed order)}.
\]
And a permutation's transpose is its inverse: $\PP^{\tp}=\PP^{-1}$ (swapping, then swapping back, is $I$).
\end{notesbox}

\boxguard
\begin{practicebox}
Let $A=\begin{bmatrix} 0 & 4 \\ 2 & 6 \end{bmatrix}$ and $\bb=\begin{bmatrix} 4 \\ 8 \end{bmatrix}$.
\begin{enumerate}[leftmargin=1.6em]
  \item Write the swap matrix $\PP$ and compute $\PP A$. \hfill $\PP A=\begin{bmatrix}\ {\color{keyred}\mathbf{2}} & {\color{keyred}\mathbf{6}}\ \\[2pt] {\color{keyred}\mathbf{0}} & {\color{keyred}\mathbf{4}} \end{bmatrix}$
  \item Solve $A\xx=\bb$ using $\PP A\xx=\PP\bb$. \hfill $x=\ans{$1$}$\quad $y=\ans{$1$}$
\begin{work}
  \PP\bb &= \begin{bsmallmatrix} 8\\4 \end{bsmallmatrix} \\
  4y &= 4,\ y = 1 \\
  2x+6 &= 8,\ x = 1
\end{work}
  \item Write $A^{\tp}$. Is $A$ symmetric?
\begin{work}
  A^{\tp} &= \begin{bsmallmatrix} 0&2\\4&6 \end{bsmallmatrix} \ne A
\end{work}
        \ansline{Not symmetric --- the off-diagonal entries $4$ and $2$ differ.}
\end{enumerate}
\end{practicebox}

\end{document}
